\section{FRFP for Organizations: Business and Governance Overview}
\label{app:business-frfp}

This appendix presents a practitioner-oriented version of the Foundational Reasoning and Feedback Protocol (FRFP).
The main body of the paper develops FRFP as a formal architecture for long-horizon human--AI scientific workflows, with explicit/tacit categories, a boundary morphism, and dynamic and institutional semantics.
Here, we restate the core ideas in non-technical terms and sketch how they can guide concrete organizational practice.

The intended audience for this appendix includes:
senior technical leaders (e.g.\ Chief Scientist, Head of AI),
risk and governance teams,
and product owners in domains where AI systems support high-stakes or long-horizon workflows.

\subsection{Executive Overview}

Organizations increasingly rely on AI systems---especially large language models and tool-using assistants---for scientific research, product development, financial analysis, and clinical or regulatory decision support.
These systems can greatly accelerate work, but they also create subtle new failure modes:
over long workflows, humans may become over-reliant on AI outputs, tacit grounding can degrade, and institutions can lose clear lines of responsibility.

FRFP offers three core ideas for structuring such workflows:

\begin{enumerate}
  \item \textbf{Explicit vs.\ tacit.}
  We distinguish between \emph{explicit} artifacts (documents, code, data, prompts, model outputs) and \emph{tacit} human understanding (expertise, contextual judgment, the sense that something is ``off'').
  AI systems operate only on explicit artifacts; correctness lives in tacit human states.

  \item \textbf{A single actionable boundary.}
  The key architectural feature is a controlled boundary between explicit and tacit spaces.
  Crossing this boundary corresponds to a human \emph{accepting} an explicit artifact as ``good enough to act on'' in some context (signing a report, approving a decision, submitting a paper).
  FRFP insists that AI systems do not cross this boundary: they may propose and transform explicit artifacts, but they do not have correctness authority.

  \item \textbf{Safe horizons instead of zero hallucination.}
  Over long horizons, humans who lean heavily on AI tend to check less, remember less, and rely more on the system.
  FRFP models this as degradation of tacit grounding and shows that, under realistic conditions, hallucinations (accepting ungrounded outputs) are inevitable along sufficiently long runs.
  The goal is therefore not ``no hallucinations ever'', but carefully designed \emph{safe horizons}: limits on how far a workflow can run before deeper human review, re-grounding, or escalation is required.
\end{enumerate}

From an organizational perspective, FRFP yields:
(i) architectural clarity about who may do what,
(ii) governance constraints and impossibility results that explain which roles cannot be fully automated,
and (iii) concrete patterns and checklists for designing AI-supported workflows that remain auditable and governable as capabilities scale.

\subsection{FRFP Architecture in Organizational Terms}

\paragraph{Explicit work.}
Explicit work consists of everything that can be written down, stored, and processed by machines:
prompts, model outputs, code, data tables, dashboards, reports, and so on.
This is where AI systems operate most naturally: generating, transforming, and organizing explicit artifacts.

\paragraph{Tacit work.}
Tacit work consists of human understanding and judgment:
domain expertise, sense of context, knowledge of institutional constraints, and the intuitive ability to say ``this looks wrong''.
Correctness---in the sense relevant for scientific, financial, or clinical responsibility---is treated as a property of tacit states, not of explicit artifacts alone.

\paragraph{The boundary.}
The FRFP boundary is the moment when a human accepts or acts on explicit material influenced by AI:
signing a clinical note, approving a risk position, submitting an AI-assisted manuscript, or green-lighting a model deployment.
At these boundary points, FRFP requires that:
(i) the responsible human role is clearly identified,
(ii) the information available to that role is specified,
and (iii) correctness is understood as their responsibility, not the system's.

AI tools may contribute to everything that happens \emph{before} the boundary (e.g.\ summaries, scenarios, drafts, analyses), but do not themselves perform the boundary crossing.

\paragraph{Roles and responsibilities.}
In an FRFP-aligned organization, it is helpful to distinguish four functional roles:

\begin{itemize}
  \item \emph{Producers of explicit artifacts:}
  teams using AI tools to generate and transform documents, models, analyses, or proposals.

  \item \emph{Tacit owners / decision-makers:}
  individuals who formally accept or reject explicit artifacts at boundary points, and who own correctness for those decisions.

  \item \emph{Institutional aggregators:}
  committees, review boards, and governance bodies that combine explicit artifacts and tacit judgments across many individuals (e.g.\ model risk committees, clinical guideline committees, internal research review boards).

  \item \emph{AI and tooling stewards:}
  platform teams responsible for AI integration, logging, and ensuring that tools respect explicit/tacit boundaries and support auditability.
\end{itemize}

\paragraph{Safe horizons.}
Every AI-assisted workflow has a safe operating horizon: a limit on how far, or how long, the workflow can run before tacit grounding and error-detection degrade unacceptably.
In practice, organizations should specify, for high-stakes processes:

\begin{itemize}
  \item how many AI-assisted steps may occur before a deeper human review is mandatory;
  \item how often humans must re-ground themselves in primary data or source documents;
  \item what kinds of anomalies or thresholds trigger escalation.
\end{itemize}

FRFP shows that, under broad qualitative assumptions, infinite safe horizons are not achievable; instead, horizon management becomes a core design responsibility.

\subsection{Examples of FRFP-Aligned Workflows}

\paragraph{Scientific research and R\&D.}
In a research workflow, AI systems may assist with literature search, clustering, summarization, hypothesis generation, and drafting.
The explicit layer includes retrieved papers, AI-generated topic maps and summaries, and draft text.
Tacit work includes: judging whether the literature coverage is adequate, assessing plausibility of hypotheses, choosing experiments, and interpreting results.

FRFP suggests:
(i) defining boundary points such as pre-registration of hypotheses and final manuscript approval, where tacit owners (e.g.\ PIs) explicitly accept responsibility;
(ii) instituting checkpoints where researchers must directly consult primary sources rather than relying solely on AI summaries;
and (iii) logging key AI-assisted reasoning steps so that the intellectual provenance of conclusions can be reconstructed.

\paragraph{Regulated financial or clinical decisions.}
In finance, AI might generate stress-test scenarios, synthesize risk reports, or help draft credit memos.
In healthcare, it might summarize charts, suggest diagnoses, or propose treatment options.
In both cases, FRFP confines AI to the generation and transformation of explicit artifacts, while reserving acceptance of credit decisions, trading positions, or clinical orders to identified human roles.

Safe horizons may be expressed as limits on:
the number of successive AI-assisted reports before independent human re-baselining,
the extent to which AI-drafted text can be reused without full review,
or the duration for which an AI-supported monitoring process may run without human reassessment of assumptions.

\paragraph{Internal governance and oversight.}
Model governance, ethics boards, and risk committees increasingly use AI to summarize documentation and surface issues.
FRFP treats governance itself as a pipeline over explicit artifacts (model cards, evaluations, incident logs, AI-generated summaries) and tacit judgments (legal, ethical, and strategic interpretation by committee members).
Key governance actions---model approval, acceptance of residual risk, adoption of policies---are boundary crossings that must remain human-exclusive.

FRFP's institutional layer implies that fully automated governance (e.g.\ ``deploy if scores exceed threshold and no alerts fire'') cannot generally provide the same correctness guarantees as human governance, even with sophisticated monitoring.
AI may and should support governance, but not replace it.

\subsection{Implementation Roadmap}

From an implementation perspective, organizations can adopt FRFP in stages:

\begin{enumerate}
  \item \textbf{Assess and map.}
  Inventory where AI is currently used; identify existing (possibly implicit) boundary points and areas with no clear boundary.

  \item \textbf{Redesign key workflows.}
  For a small number of high-impact workflows, explicitly map explicit vs.\ tacit steps, define boundary roles and information requirements, and choose safe horizons and checkpoints.
  Decide what should be logged to reconstruct AI-assisted reasoning.

  \item \textbf{Pilot and refine.}
  Run FRFP-aligned versions of these workflows with motivated teams.
  Monitor incidents, overrides, and user experience; adjust boundaries, horizons, and logging based on what works in practice.

  \item \textbf{Standardize and scale.}
  Convert successful pilots into reusable FRFP templates (for research, risk decisions, clinical documentation, model governance, etc.).
  Integrate FRFP design into product lifecycles, model approval processes, and change management.

  \item \textbf{Embed in governance and training.}
  Update AI use policies, model governance frameworks, and risk documents to reflect explicit/tacit boundaries, human-exclusive correctness, and limits on automation.
  Train practitioners and leaders in FRFP concepts and their responsibilities at boundary points.
\end{enumerate}

\subsection{Common Pitfalls and How FRFP Helps Avoid Them}

In applying FRFP, several recurring pitfalls are worth highlighting:

\begin{itemize}
  \item \textbf{Over-automation disguised as decision support.}
  Tools described as ``assistants'' or ``copilots'' effectively become decision-makers when humans never override them.
  FRFP reacts by treating behavior, not labels, as the signal: if a system is de facto making decisions, boundaries must be redesigned.

  \item \textbf{Boundary erosion over time.}
  Initial review steps can become rubber-stamping under time pressure.
  FRFP recommends making boundary decisions observable (e.g.\ by logging changes and overrides) and periodically auditing samples to detect erosion.

  \item \textbf{Metric overconfidence.}
  Organizations may over-trust dashboards, scores, or evaluation metrics.
  Under FRFP, such metrics are inputs to tacit judgment, not substitutes; governance documents should explicitly state what metrics do \emph{not} capture.

  \item \textbf{Checklist theatre.}
  Long, low-value checklists can create a false sense of safety.
  FRFP encourages short, targeted prompts that trigger real reflection (e.g.\ ``What primary sources did you personally review?''), and regular pruning of items that do not change behavior.

  \item \textbf{Treating FRFP as pure compliance.}
  If FRFP is seen only as an external requirement, teams will optimize for appearance.
  Embedding FRFP in product and research strategy---not just policy---helps keep the focus on actual correctness and robustness.
\end{itemize}

\subsection{FRFP Design Checklist}

For ease of adoption, we collect here a short checklist that teams can use when designing or reviewing AI-supported workflows.

\begin{itemize}
  \item \textbf{Workflow mapping:}
    have we charted the main steps and labeled each as explicit or tacit?

  \item \textbf{Boundaries:}
    where exactly are the decision points where correctness matters,
    who is responsible at each, and what information must they see?

  \item \textbf{AI's role:}
    are AI systems confined to explicit work (generation, transformation, organization),
    or are they implicitly making decisions that should be human-held?

  \item \textbf{Safe horizons:}
    for this workflow, how far can AI-assisted reasoning proceed before requiring deeper human review, re-grounding, or escalation?
    what are the triggers for resets?

  \item \textbf{Logging and traceability:}
    are prompts, key AI outputs, and boundary decisions logged in a way that lets us reconstruct how a conclusion was reached?

  \item \textbf{Governance:}
    which committee or body oversees this workflow?
    do they understand the explicit/tacit boundaries and have the authority to adjust them?

  \item \textbf{Training and culture:}
    have the people using the workflow been trained in their responsibilities at boundary points and encouraged to challenge AI outputs when needed?
\end{itemize}

Used consistently, this checklist can help ensure that FRFP is not just a theoretical architecture but a practical guide to designing AI-supported work that is safer, more auditable, and more robust against long-horizon drift.
